\section{Representations of matrix algebras and coalgebras}

\begin{lemma}[Splitting an arbitrary idempotent]
    \label{thm:glm_boundary_idempotent}
    \lean{Matrix.exists_rankFactorization_of_idempotent}
    \leanok
    Every idempotent $P\in M_D(K)$ over a field factors as $P=WV$,
    $VW=I_m$, with $m=\operatorname{rank}P$. This includes $m=0$ and
    imposes no self-adjointness condition.
\end{lemma}

\begin{proof}
    \leanok
    Choose a basis of $\operatorname{ran}P$. Let $W$ include these
    coordinates in $K^D$, and let $V$ take coordinates after applying $P$.
    The identity $P^2=P$ says that $P$ is the identity on its range, hence
    $VW=I_m$ and $WV=P$.
\end{proof}

\begin{lemma}[Biorthogonal matrix-algebra representations]
    \label{thm:glm_boundary_matrix_blocks}
    \lean{Matrix.exists_matrixUnitFactorization,
        Matrix.sum_multiplicityBlocks_mul,
        Matrix.exists_matrixUnitBlocks,
        Matrix.exists_multiplicativeLinearMap_blocks,
        Matrix.exists_multiplicativeLinearMap_blocks_of_fintype,
        Matrix.biorthogonal_sum_support,
        Matrix.exists_piMatrix_blocks}
    \leanok
    A multiplicative $K$-linear map $\rho\colon\prod_c M_{D_c}(K)\to M_D(K)$
    admits finite multiplicities $m_c$ and matrices
    $W_{c\mu}\colon K^{D_c}\to K^D$, $V_{c\mu}\colon K^D\to K^{D_c}$ such that
    \begin{align}
        V_{c\mu}W_{d\nu}
                & =\delta_{cd}\delta_{\mu\nu}I_{D_c},
        \label{eq:glm_boundary_matrix_orthogonal}       \\
        \rho(M) & =\sum_{c,\mu}W_{c\mu}M_cV_{c\mu},   &
        \rho(1) & =\sum_{c,\mu}W_{c\mu}V_{c\mu}.
        \label{eq:glm_boundary_matrix_reconstruct}
    \end{align}
    Zero-dimensional factors and zero multiplicities are allowed. This is
    the decomposition of \cite[Appendix~A]{GarreRubio2023MPOSym}.
\end{lemma}

\begin{proof}
    \leanok
    \uses{thm:glm_boundary_idempotent}
    For a nonzero factor, put $E_{ij}=\rho(e_{c,ij})$ and choose
    $E_{00}=CH$ with $HC=I$. The identities
    $E_{ij}E_{kl}=\delta_{jk}E_{il}$ give
    $(HE_{0i})(E_{j0}C)=\delta_{ij}I$ and
    $(E_{i0}C)(HE_{0j})=E_{ij}$. Separate the common coordinate into
    multiplicity indices and expand an arbitrary matrix in matrix units.
    Distinct factors have annihilating support idempotents, so their
    rectangular factors are mutually biorthogonal. Summing the factor
    expansions gives (\ref{eq:glm_boundary_matrix_reconstruct}).
\end{proof}

\begin{definition}[Pre-bialgebra]
    \label{def:mpoalg_pre_bialgebra}
    \lean{PreBialgebra, Bialgebra.toPreBialgebra}
    \leanok
    A pre-bialgebra over a commutative semiring $R$ is an algebra
    $\mathcal A$ over $R$ with linear maps
    $\Delta\colon\mathcal A\to\mathcal A\otimes\mathcal A$ and
    $\epsilon\colon\mathcal A\to R$ such that $\Delta$ is multiplicative and
    coassociative and $\epsilon$ is a counit:
    \begin{align}
        \Delta(xy)                    & =\Delta(x)\Delta(y),            &
        (\Delta\otimes\Id)\circ\Delta & =(\Id\otimes\Delta)\circ\Delta,
        \label{eq:mpoalg_pre_bialgebra}                                   \\
        (\epsilon\otimes\Id)\circ\Delta
                                      & =\Id,                           &
        (\Id\otimes\epsilon)\circ\Delta
                                      & =\Id,
        \label{eq:mpoalg_pre_bialgebra_counit}
    \end{align}
    after the standard identifications of the two triple tensor products and
    of $R\otimes\mathcal A$ and $\mathcal A\otimes R$ with $\mathcal A$.
    Neither $\Delta(1)=1\otimes1$ nor any relation between $\epsilon$
    and the unit or the product is required. Every bialgebra is a
    pre-bialgebra.

    This is \cite[Definition~4.1]{Molnar2022MPOAlgebrasI}: an algebra that is
    also a coalgebra \cite[Definition~3.1]{Molnar2022MPOAlgebrasI} with
    multiplicative coproduct. The source works with finite-dimensional
    complex vector spaces.
\end{definition}

\begin{lemma}[Same-order involution of the convolution dual]
    \label{thm:glm_coalgebra_dual_star}
    \lean{Coalgebra.dualStar_mul, Coalgebra.dualStarMulEquiv,
        Coalgebra.dualStar_one}
    \leanok
    Let $H$ be a coalgebra with an involution such that
    $\Delta(x^*)=\Delta(x)^*$. For $f\in H^*$, define
    $\kappa(f)(x)=\overline{f(x^*)}$. Then $\kappa$ is conjugate-linear,
    involutive, and satisfies
    \begin{align}
        \kappa(fg) & =\kappa(f)\kappa(g), \notag \\
        \kappa(1)  & =1.
        \label{eq:glm_dual_star_product}
    \end{align}
    Here multiplication is convolution, and $1$ is the counit.
\end{lemma}

\begin{proof}
    \leanok
    On a pure tensor, complex conjugation sends
    $f(x^*)g(y^*)$ to $\kappa(f)(x)\kappa(g)(y)$.
    Extend this identity additively to $\Delta(x)$ and use the
    coproduct compatibility. Involutivity makes $\kappa$ a bijective
    multiplicative map, so it fixes the unit.
\end{proof}

\begin{lemma}[Representations of the Klein four-group by involutions]
    \label{lem:elementary_klein_involutions}
    \lean{Matrix.ofCommutingInvolutions,
        Matrix.ofCommutingInvolutions_mul_conjTranspose,
        Matrix.mul_of_anticommuting_involutions,
        Matrix.mul_of_anticommuting_neg_involution}
    \leanok
    Let $P_1,P_2$ be square complex matrices with $P_1^2=P_2^2=\Id$. If
    $P_1P_2=P_2P_1$, then $(g_1,g_2)\mapsto P_1^{g_1}P_2^{g_2}$ is a
    representation of $\mathbb Z_2\times\mathbb Z_2$, unitary when $P_1$ and
    $P_2$ are. If instead $P_2P_1=-P_1P_2$, the same assignment satisfies
    $u(g)u(h)=(-1)^{g_2h_1}u(gh)$; if $P_1^2=-\Id$ and $P_2P_1=-P_1P_2$, it
    satisfies $u(g)u(h)=(-1)^{(g_1+g_2)h_1}u(gh)$.
\end{lemma}

\begin{proof}
    \leanok
    Check the sixteen pairs of group elements using the involution and
    commutation relations.
\end{proof}
